\documentclass[11pt,a4paper]{article}
\usepackage[margin=2.2cm]{geometry}
\usepackage[T1]{fontenc}
\usepackage{lmodern}
\usepackage{booktabs,tabularx,array,enumitem}
\usepackage[hidelinks]{hyperref}
\usepackage{parskip}
\newcolumntype{Y}{>{\raggedright\arraybackslash}X}
\setlist{itemsep=4pt,topsep=4pt}
\begin{document}

\begin{center}
{\LARGE\bfseries Phone Test: Easy Steps}\\[4pt]
{\large No cable, no router. Just a phone hotspot and this laptop.}
\end{center}

\section*{What you need}
\begin{itemize}
\item This laptop.
\item One phone to make a hotspot (yours).
\item The phones to test: yours and friends' (ask them first). Best to have: an \textbf{iPhone with Safari}, an \textbf{Android with Chrome}, and \textbf{Firefox} on Android if possible.
\item About 2 minutes per phone.
\end{itemize}

\section*{Part 1: Start (once)}
\begin{enumerate}
\item On your phone, turn on \textbf{Mobile Hotspot}. Mobile data must be on.
\item Connect the \textbf{laptop's Wi-Fi} to that hotspot.
\item On the laptop, open the folder \texttt{Desktop\textbackslash cnfinalpaper\textbackslash hardware\_test}.
\item Click the folder's address bar at the top, type \texttt{cmd}, press \textbf{Enter}. A black window opens.
\item Type \texttt{python server.py} and press \textbf{Enter}.
\item It prints one or more addresses like \texttt{http://192.168.43.12:8000/}. Write down the one that starts with \texttt{192.168}, \texttt{172.20} or \texttt{10.}
\item If Windows asks about the firewall, tick \textbf{both} Private and Public, then press \textbf{Allow}.
\item Keep the black window open until you finish.
\end{enumerate}

\section*{Part 2: Test each phone}
\begin{enumerate}
\item Connect the phone to the \textbf{same hotspot}. (The hotspot phone itself can also be tested; if the page will not open on it, skip it.)
\item Open the browser you want to test (Safari, Chrome or Firefox).
\item Type the address from Part 1, step 6, exactly, including \texttt{:8000}. A page called \textbf{WebRTC path test} opens.
\item In the box \textbf{This device and browser}, type the phone and browser, for example \texttt{iPhone 13 Safari} or \texttt{Redmi Note 12 Chrome}.
\item In the box \textbf{Network setup}, type \texttt{hotspot}.
\item Press \textbf{Test this device only}. Wait about 10 seconds.
\item The screen shows \textbf{Result: HIDES local address} or \textbf{Result: REVEALS local address}. It is saved on the laptop automatically.
\item \textbf{Reload the page and press the button again}, so each phone is tested \textbf{3 times}.
\item Repeat for every phone and every browser.
\end{enumerate}

\section*{Part 3: Write down (once per phone)}
\begin{tabularx}{\textwidth}{@{}YYYY@{}}
\toprule
\textbf{Phone model} & \textbf{OS version} (Settings $\rightarrow$ About) & \textbf{Browser and version} (browser $\rightarrow$ About) & \textbf{Result (3 runs)} \\
\midrule
 & & & \\[14pt]
 & & & \\[14pt]
 & & & \\[14pt]
 & & & \\[14pt]
\bottomrule
\end{tabularx}

\section*{Part 4: Finish}
\begin{enumerate}
\item In the black window press \textbf{Ctrl + C}, then close it.
\item A file \texttt{results.jsonl} is now in the \texttt{hardware\_test} folder. It holds only address \emph{types}, never the actual addresses.
\item Send that file and the table above to Claude to add to the paper.
\end{enumerate}

\section*{If something goes wrong}
\begin{tabularx}{\textwidth}{@{}p{5.2cm}Y@{}}
\toprule
\textbf{Problem} & \textbf{Fix} \\
\midrule
Page does not open on the phone & Check the phone is on the same hotspot. Try the other address the black window printed. Allow Python in the firewall for Public networks. \\
\texttt{python} is not recognised & Close the black window and open it again from the folder; if it still fails, ask Claude. \\
Result: NO CANDIDATES & Turn mobile data on; the hotspot needs internet. \\
iPhone hotspot: laptop cannot connect & On the iPhone, keep the Personal Hotspot screen open while connecting. \\
\bottomrule
\end{tabularx}

\vspace{6pt}
\textbf{What this test does and does not show.} It shows what each phone browser reveals about its own address. It does not test client isolation, because a phone hotspot has no such setting. Each phone talks only to the laptop and to Google's public STUN server, as in a normal video call.

\end{document}
